\chapter{集合论复习}

\section{集合}

下面集合的有关概念和其运算性质.

\begin{definition}[幂集][def:1.1]\index{m@幂集}
    设$X$是一个集合，$X$的所有子集组成的集合称为$X$的\colorbf{幂集}，记作$2^X$.
\end{definition}

这个记号的一个合理解释是，若一个集合有$X$个元素，那么它有$2^X$个子集.

\begin{definition}[集合的卡氏积][def:1.2]\index{k@卡氏积}
    设$X_1,X_2,\cdots,X_n$为集合，则它们的\colorbf{卡氏积}定义为$X_1\times X_2\times\cdots\times X_n := \big\{(x_1,x_2,\cdots,x_n)|x_i\in X_i\big\}$.
\end{definition}

\begin{proposition}[集合运算的分配律][prop:1.1]
    \wideline
    \begin{enumerate}
        \item $A\cup \left(\bigcap_{\lambda\in\Lambda}B_\lambda\right)=\bigcap_{\lambda\in\Lambda}{\left(A\cup B_\lambda\right)}$
        \item $A\cap \left(\bigcup_{\lambda\in\Lambda}B_\lambda\right)=\bigcup_{\lambda\in\Lambda}{\left(A\cap B_\lambda\right)}$
        \item $A\times \left(\bigcup_{\lambda\in\Lambda}B_\lambda\right)=\bigcup_{\lambda\in\Lambda}{\left(A\times B_\lambda\right)}$
        \item $A\times \left(\bigcap_{\lambda\in\Lambda}B_\lambda\right)=\bigcap_{\lambda\in\Lambda}{\left(A\times B_\lambda\right)}$
        \item $A\times \left(B_1\setminus B_2\right)=(A\times B_1)\setminus (A\times B_2)$
    \end{enumerate}
\end{proposition}
上面的$\Lambda$是一个指标集，它意味着上述集合的性质是否满足和集合的数量无关，也就是哪怕有不可数个集合，它们也一定满足上述的运算性质.

\begin{theorem}[德摩根定律][thm:1.1]\index{d@德摩根定律}
    \wideline
    \begin{enumerate}
        \item $A\setminus \left(\bigcup_{\lambda\in\Lambda}{B_{\lambda}}\right)=\bigcap_{\lambda\in\Lambda}{\left(A\setminus B_{\lambda}\right)}$
        \item $A\setminus \left(\bigcap_{\lambda\in\Lambda}{B_{\lambda}}\right)=\bigcup_{\lambda\in\Lambda}{\left(A\setminus B_{\lambda}\right)}$
    \end{enumerate}
\end{theorem}

\section{映射}

下面给出映射的定义和相关概念，以及映射的一些性质.

\begin{definition}[映射][def:1.3]\index{y@映射}
    设$X,Y$是两个集合，则$X,Y$之间的\colorbf{映射} $f:X\to Y$将$X$中的每个元素$x$ \uline{唯一地}对应到$Y$中一个元素$y$.其中$X$称为\colorbf{定义域}，$Y$称为\colorbf{值域}.
\end{definition}

\begin{definition}[映射的像][def:1.4]\index{x@像}
    设$f:X\to Y$是一个映射，则子集$A\subseteq X$在$f$之下的\colorbf{像}为
    \[f(A)=\big\{y\in Y\big|\exists\, x\in A,\,\st f(x)=y \big\}\subseteq Y\]
\end{definition}

\begin{definition}[原像][def:1.5]\index{y@原像}
    设$f:X\to Y$是一个映射，则子集$B\subseteq Y$在$f$之下的\colorbf{原像}为
    \[f^{-1}(B)=\big\{x\in X\big|f(x)\in B \big\}\subseteq X\]
\end{definition}

\begin{attention}
    注意，这里的$f^{-1}$不表示$f$的逆映射或是$f$的反函数.
\end{attention}

\begin{definition}[单射,满射,双射][def:1.6]\index{d@单射}\index{m@满射}\index{s@双射}
    设$f:X\to Y$是一个映射，如果：
    \begin{enumerate}
        \item $\forall\, x,y\in X,f(x)=f(y)\Rightarrow x=y$，则称$f$是一个\colorbf{单射}.
        \item $\forall\,y\in Y,\exists\, x\in X,\st f(x)=y$，则称$f$是一个\colorbf{满射}.
        \item 若$f$即是单射也是满射，则称$f$是一个\colorbf{双射}.
    \end{enumerate}
\end{definition}
对于单射，我们通常还有其等价条件，即如果$\forall\, x,y\in X,x\neq y\Rightarrow f(x)\neq f(y)$，则称$f$是一个单射.这个条件和上面的条件互为逆否命题，因此可以很容易知道两个条件实质上是完全相同的.





\begin{definition}[逆映射][def:1.7]\index{n@逆映射}
    如果$f:X\to Y$是一个双射，定义$f^{-1}:Y\to X$是$f$的\colorbf{逆映射}，满足
    \begin{equation*}
    f:x\mapsto y \Rightarrow f^{-1}:y\mapsto x.
    \end{equation*}
\end{definition}

\begin{proposition}[映射的性质][prop:1.2]
    \wideline
    如果$f:X\to Y$，$g:Y\to Z$是两个映射，$A_\lambda\subset X,B_\lambda \subset Y$，则
    \begin{enumerate}
        \item $f^{-1}\left(\bigcup_{\lambda}{B_\lambda}\right)=\bigcup_{\lambda}{f^{-1}(B_\lambda)}$
        \item $f^{-1}\left(\bigcap_{\lambda}{B_\lambda}\right)=\bigcap_{\lambda}{f^{-1}(B_\lambda)}$
        \item $f\left(\bigcup_{\lambda}{B_\lambda}\right)=\bigcup_{\lambda}{f(B_\lambda)}$
        \item $f\left(\bigcap_{\lambda}{B_\lambda}\right)\subset\bigcap_{\lambda}{f(B_\lambda)}$，当$f$为单射时，等号成立.
        \item $f\left(f^{-1}(B)\right)\subset B$，当$f$为满射时，等号成立.
        \item $f^{-1}\left(f(A)\right)\supset A$，当$f$为单射时，等号成立.
        \item $(g\circ f)(A)=g \left(f(A)\right),A\subset X$
        \item $(g\circ f)^{-1}(C)=f^{-1}\left(g^{-1}(C)\right),C\subset Z$
    \end{enumerate}
\end{proposition}
\begin{proof}
    \begin{enumerate}
        \setcounter{enumi}{3}
        \item $\forall\, x\in \bigcap_{\lambda}{B_\lambda}$，有$\forall\,\lambda\in\Lambda,x\in B_{\lambda} \Rightarrow \forall\,\lambda\in\Lambda,f(x)\in f \left(B_\lambda\right)$\par
        因此$\forall\,\lambda\in\Lambda,f\left(\bigcap_{\lambda}{B_\lambda}\right)\subset f(B_\lambda)\Rightarrow f\left(\bigcap_{\lambda}{B_\lambda}\right)\subset\bigcap_{\lambda}{f(B_\lambda)}$.\par
        当$f$为单射时，对任意$y\in \bigcap_{\lambda}{f(B_\lambda)}$，都有$y\in f(B_\lambda)(\forall \lambda)$，所以$\forall \lambda,\exists\, x_\lambda\in B_\lambda,\st f(x_\lambda)=y$.\par
        由于$f$为单射，可知所有的$x_\lambda$都相等，记为$x_0$，即$\forall \lambda,x_\lambda=x_0 \in X$，则$\forall \lambda,x_0\in B_\lambda\ \Rightarrow x_0\in\bigcap_{\lambda}{B_\lambda} \Rightarrow f(x_0)=y\in f\left(\bigcap_\lambda{B_\lambda}\right) \Rightarrow \bigcap_{\lambda}{f(B_\lambda)}\subset f\left(\bigcap_\lambda{B_\lambda}\right)$.\par
        综上可知当$f$为单射时，$f\left(\bigcap_{\lambda}{B_\lambda}\right) = \bigcap_{\lambda}{f(B_\lambda)}$
        \item $\forall\, x\in f^{-1}(B)$，$\exists\, y\in B,\st f(x)=y\in B \Rightarrow f\left(f^{-1}(B)\right)\subset B$.\par
        当$f$为满射时，$\forall\, y\in B,\exists\, x\in f^{-1}(B),\st f(x)=y \Rightarrow B \subset f(f^{-1}(B))$，因此$f\left(f^{-1}(B)\right)= B$.
        \item $\forall\, x\in A,\exists\, y\in f(A),\st f(x)=y$，故$x\in f^{-1}(y) = f^{-1}(f(x)) \subset f(f^{-1}(A)) \Rightarrow A\subset f^{-1}(f(A))$.\par
        当$f$为单射时，$\forall y\in f(A),\exists!\, x\in A,\st f(x)=y$，故$f^{-1}(y)=x\in A \Rightarrow f^{-1}(f(A))\subset A$.\par
        因此，当$f$为单射时，$f^{-1}(f(A))=A$.\qedhere
    \end{enumerate}
\end{proof}

上述是4-6条映射性质的证明，其余性质较为简单故不作证明，所使用证明两个集合相等的方式是证明两个集合互相包含.


\begin{definition}[映射的限制][def:1.8]
    设$f:X\to Y$是一个映射，$A\subset X$，定义$f$在$A$上的\colorbf{限制}为
    \begin{align*}
        f|_A :A &\to F \\
        a &\mapsto f(a)
    \end{align*}
    反过来，$f$称为$f|_A$在$X$上的\colorbf{扩张}.
\end{definition}

\begin{definition}[特殊的映射][def:1.9]
    \begin{enumerate}
        \item ${\setjot[-2]\begin{aligned}[t]
            \mathrm{id}_X: X &\to X \\
            x &\mapsto x
        \end{aligned}}$称为$X$上的\colorbf{恒等映射}.
        \item 若$A\subset X$，则${\setjot[-2]\begin{aligned}[t]
            i:A &\to X \\
            a &\mapsto a
        \end{aligned}}$称为$A$到$X$的\colorbf{包含映射}.
        \item $\setjot[-2]\begin{aligned}[t]
            \Delta_X:X &\to X\times X \\
            x &\mapsto (x,x)
        \end{aligned}$称为$X$的\colorbf{对角映射}.
    \end{enumerate}
\end{definition}

\begin{remark}
    关于包含映射，我们有以下交换图： \par
    \begin{wrapfigure}{r}{0.25\textwidth}
        \vspace{-4 em}
        \begin{tikzcd}[labels={font=\normalsize}]
            A \arrow[rr, "f|_A"] \arrow[rd, "i"'] &                    & Y \\
                                      & X \arrow[ru, "f"'] &  
        \end{tikzcd}
        \vspace{-2 em}
    \end{wrapfigure}
    \par

也就是说，我们可以把限制映射看成包含映射和映射的复合，即$f|_A=f\circ i$.

\end{remark}

\section{等价关系}

\begin{definition}[等价关系][def:1.10]\index{d@等价关系}
    设$\sim$是集合$X$上的一个二元关系，如果它满足以下条件：
    \begin{enumerate}
        \item 自反性：$\forall\, x\in X,\,x\sim x$
        \item 对称性: $\forall\, x,y\in X,\,x\sim y\Rightarrow y\sim x$
        \item 传递性: $\forall\, x,y,z\in X,\,x\sim y\,\text{且}\,y\sim z\Rightarrow x\sim z$
    \end{enumerate}
    则称$\sim$为$X$上的一个\colorbf{等价关系}.
\end{definition}



\begin{definition}[等价类][def:1.11]\index{d@等价类}
    设$\sim$是集合$X$上的一个等价关系，则$A$中所有等价于$x$的元素组成的集合称为$x$的\colorbf{等价类}，记作$[x]=\{y\in X|x\sim y\}$.
\end{definition}
我们经常把一个等价类$[x]$中的一个元素$x$称作这个等价类的\bluebf{代表元}.

\begin{proposition}[等价类的性质][prop:1.3]
    两个等价类要么是相等的，要么是不交的.
\end{proposition}
\begin{proof}
    使用反证法，假设$[x]$和$[y]$是集合$X$的两个不相等的等价类且$[x]\cap[y]\neq\varnothing$，则$\exists\, z\in [x]\cap[y]$.\par
    由等价类的定义可知$\forall\, \tilde{x}\in[x],\tilde{x}\sim z$.同理$\forall\, \tilde{y}\in[y],\tilde{y}\sim z$.\par
    $\xLongrightarrow{\text{传递性}}\forall\, \tilde{x}\in[x],\forall\,\tilde{y}\in[y],\tilde{x}\sim\tilde{y} \Rightarrow [x]=[y]$，矛盾!.\par
    因此$[x]\cap[y]=\varnothing$，即若$[x]$和$[y]$不相等，它们必定不交.
\end{proof}

这个性质实际上告诉我们等价类可以把一个集合完整地划分为不同的小集合.


\begin{definition}[商集][def:1.12]\index{s@商集}
    设$\sim$是集合$X$上的一个等价关系，所有等价类的集合称为\colorbf{商集}，记为$X/\sim \;= \left\{[x]|x\in X\right\}$.
\end{definition}
从元素的角度来看，从一个集合的所有等价类中各选取一个代表元组成的集合就是它的商集.


